\section{Two slopes suffice at a scalar interface}
\label{sec:messages}

Consider a subtree $T_v$, with parent interface $\tau$. Write $f_v$ for the sum
of objective contributions of its private coordinates. A compatible assignment
means one satisfying all equations, nonnegativity conditions and forbidden
pairs inside the subtree.

\begin{definition}[Continuation message]
The zero-port flag $Z_v$ is true when a compatible assignment satisfies
$\tau=0$ and $f_v>0$. For $s\in\{-1,1\}$, sign $s$ is available when some
compatible assignment has $\tau=s$. When $Z_v$ is false and sign $s$ is available,
put
\[
 \alpha_{v,s}=\max\{f_v:\tau=s,\ \text{the subtree assignment is compatible}\}.
\]
An unavailable sign is represented by a separate symbol, not by the numerical
value zero.
\end{definition}

\begin{theorem}[Exact scalar continuation]
\label{thm:message}
If $Z_v$ is false, each available $\alpha_{v,s}$ is finite, attained, and rational.
For every specified scalar $t$, the best subtree objective is exactly
\[
 F_v(t)=
 \begin{cases}
 0,&t=0,\\
 |t|\alpha_{v,\operatorname{sgn}t},&t\ne0\text{ and its sign is available},\\
 -\infty,&t\ne0\text{ and its sign is unavailable}.
 \end{cases}
\]
If $Z_v$ is true, the full tree already has a positive assignment, obtained by
extending a positive zero-port witness by zero outside the subtree.
\end{theorem}

\begin{proof}
Choose one allowed coordinate face for every local compatibility condition and
one sign for every internal scalar. Replacing a scalar of fixed sign by a
nonnegative variable gives a rational polyhedral cone. There are finitely many
such choices. Intersecting each cone with $\tau=s$ produces a rational
polyhedron, possibly empty.

If the objective is unbounded above on a nonempty such polyhedron, linear
programming gives a recession direction with parent coordinate zero and
positive objective. The direction remains in the same coordinate face and
sign chamber, so it is compatible. This would make $Z_v$ true. Thus every
nonempty normalized polyhedron has finite optimum. A finite optimum of a
rational polyhedron is attained at a rational point. Taking the maximum over
the finitely many nonempty choices proves the first assertion.

When $t\ne0$, multiplication or division by $|t|$ is a bijection between the
fiber at $\operatorname{sgn}t$ and the fiber at $t$, preserving compatibility.
The objective scales by $|t|$, giving the displayed formula. At zero, no
positive objective is possible by hypothesis, and the zero assignment attains
zero. Finally, a zero-parent-port assignment satisfies the outside equations
when every outside coordinate and port is set to zero. Homogeneity and the
absence of cross-bag conflicts prove the last assertion.
\end{proof}

\begin{corollary}[Context equivalence for the decision query]
Replacing a subtree by its exact two-slope message preserves whether any
homogeneous context attached only at that scalar has positive total objective.
The message need not preserve the set of internal surfaces, component counts,
full normal coordinates, or arbitrary topological attachment data.
\end{corollary}

The theorem is not a convex-hull relaxation. Availability is retained for both
signs, the value at zero is treated separately, and the best witness for one
sign need not have the same compatibility pattern as the best witness for the
other. The reduction is valid because the context sees precisely one scalar
and one additive objective contribution.

\subsection{Local recurrence}

At a vertex $v$, fix signs $\sigma_e\in\{-1,1\}$ on all incident edges and
write $\tau_e=\sigma_e z_e$, $z_e\geq0$. The local homogeneous cone in
$(x_v,z)$ is
\[
 K_{v,\sigma}=\left\{(x_v,z)\geq0:
 B_vx_v+\sum_{e\ni v}\sigma_e h_{v,e}z_e=0\right\}.
\]
Unavailable child signs force their corresponding $z_e$ to zero. Caller-disabled
original coordinates also vanish. These are coordinate faces, not new general
inequalities. On this cone the effective local objective is linear:
\begin{equation}\label{eq:effective}
 g_{v,\sigma}(x_v,z)=c_v^Tx_v+
 \sum_{e\text{ child of }v}\alpha_{e,\sigma_e}z_e.
\end{equation}
Terms on forced-zero child coordinates are omitted.

Take every local extreme ray that satisfies the local forbidden pairs and
coordinate-zero restrictions. A positive value of $g_{v,\sigma}$ on a ray with
zero parent port produces $Z_v$. Otherwise, for a ray with parent value $z_p>0$,
its candidate parent slope is $g_{v,\sigma}/z_p$. Maximize these candidates
separately for the two parent signs. At the root there is no parent port;
a positive effective value answers the original query.

\begin{proposition}[Completeness of the local-ray recurrence]
\label{prop:recurrence}
The recurrence computes the messages of Theorem~\ref{thm:message} exactly.
It remains exact when prepared ray lists are reused after any further original
coordinate deletions and any change of the linear objective.
\end{proposition}

\begin{proof}
In a fixed sign chamber, a compatible vector lies in a coordinate face on which
its conflicting coordinates vanish. The extreme rays of this face are extreme
rays of the unrestricted local cone and have support contained in that of the
vector. They therefore satisfy all its incompatibility and zero conditions.

Express the vector as a nonnegative combination of those rays. If its parent
port is zero, every constituent ray has zero parent port, since that coordinate
is nonnegative in the chamber. Positive total objective then requires one
positive zero-port ray. If its parent port is positive and no positive
zero-port ray exists, zero-port terms cannot improve the objective. The ratio
of objective to parent port on the remaining sum is a weighted average of the
ratios on its positive-port rays. Its maximum is therefore attained by one ray.

By induction, a child slope is attainable by a child assignment with the exact
required sign and scale. Independently attaching those child assignments
realizes every retained local candidate. No incompatible choice is combined
inside a private coordinate block. This proves both inequalities needed for
exactness. Further coordinate deletions select additional faces, whose rays
are obtained simply by filtering the prepared lists. Objectives change scores,
not the underlying rays.
\end{proof}

\subsection{Witness recipes and negative evidence}

A maximizing message retains one local ray and references the chosen child
recipes. To reconstruct a witness, scale a child's stored ray so that its
parent scalar equals the requested scalar. Set an entire child subtree to zero
when that scalar is zero. A zero-port positive witness found below the root is
similarly extended by zero outside its origin subtree.

The implementation uses an explicit stack, not recursion proportional to the
height of the input tree. It emits original coordinates and signed edge fluxes.
A positive checker needs only exact equation checks, nonnegativity, local pair
checks, zero-coordinate checks, and evaluation of the original objective.

For a negative answer, the supplied packet contains complete local ray lists
and all finite slopes or unavailable symbols. An independent replay enumerates
local positive supports rather than the producer's active zero sets, compares
the ray lists, and checks the scalar recurrence. This checker is exponential
in local bag size. Its existence is not a polynomial-size negative certificate
claim for arbitrary bags. Its support-enumeration cost is bounded by
$\poly(N)\sum_v2^{b_v+2\delta_v}$ with the parameters defined below; it need
not be polynomial in the smaller active-zero-set budget alone.

\section{Compiling rank-one column cuts from the source}
\label{sec:compiler}

Suppose a rooted tree carries a partition $(I_v)_v$ of the original columns of
$A$. Empty internal bags are allowed. Every forbidden pair must lie wholly in
one bag. For a nonroot vertex $v$, let $S_v$ be the union of the column bags
in its rooted subtree. Define
\begin{equation}\label{eq:cutrank}
 \lambda_v=\dim\bigl(\im A_{S_v}\cap\im A_{\overline{S_v}}\bigr)
 =\rank A_{S_v}+\rank A_{\overline{S_v}}-\rank A.
\end{equation}
The equality holds because the sum of the two column spaces is $\im A$.
The relevant quantity is this column-space intersection dimension, not the
number of geometric boundary faces, the graph cut size, or the rank of an
arbitrarily selected submatrix of crossing entries.

Assume $\lambda_v\leq1$ on every tree edge. If it equals one, choose a primitive
integer generator $w_v$ of the intersection line. If it equals zero, set
$w_v=0$ and require $\tau_v=0$ at both endpoints. For nonzero $w_v$, define the
signed flux by
\begin{equation}\label{eq:flux}
 A_{S_v}x_{S_v}=w_v\tau_v.
\end{equation}
The scalar can have either sign; replacing it by its absolute value would be
an unsound source transformation.

The compiled equation at vertex $v$ is
\begin{equation}\label{eq:node}
 A_{I_v}x_{I_v}+
 \sum_{u:\operatorname{parent}(u)=v}w_u\tau_u
 -\begin{cases}w_v\tau_v,&v\ne0,\\0,&v=0\end{cases}=0.
\end{equation}
The objective and all forbidden pairs are copied to their unique original
coordinate bags. No elimination of original coordinates is needed.

\begin{theorem}[Exact source compilation]
\label{thm:compiler}
Projection to the original coordinates is a linear bijection from the solution
space of the compiled equations to $\ker A$. The inverse is given by the unique
fluxes in \eqref{eq:flux}, with zero flux on rank-zero cuts. This bijection
preserves nonnegativity of original coordinates, every source forbidden pair,
original coordinate deletions, and every linear objective on original
coordinates. Compilation and independent source verification have polynomial
rational bit complexity in the source and supplied-tree encoding.
\end{theorem}

\begin{proof}
Given $Ax=0$, the vector $A_{S_v}x_{S_v}$ belongs to $\im A_{S_v}$ and equals
$-A_{\overline{S_v}}x_{\overline{S_v}}$, so it lies in the displayed intersection.
It therefore determines the unique stated scalar. Partitioning $S_v$ into
$I_v$ and its child subtrees yields \eqref{eq:node}.

Conversely, sum \eqref{eq:node} over the entire tree. Each nonroot flux term
appears once positively and once negatively, leaving $Ax=0$. Summing only over
a rooted subtree gives \eqref{eq:flux}. Thus no spurious scalar solutions
survive, including on edges with zero intersection. This proves bijectivity.
The additional claims follow because original coordinates and local forbidden
pairs have not been altered.

Exact Gaussian elimination computes the two column-space bases and their
intersection. Rational basis entries, intersection generators and their
primitive integer scalings have polynomial encoding length by determinant
bounds. There are polynomially many cuts and local row records. An independent
checker can compute the three ranks in \eqref{eq:cutrank}, check that a proposed
nonzero $w_v$ lies in both column spaces, require zero exactly when the
intersection has dimension zero, and compare every compiled row with
\eqref{eq:node}. These are polynomial exact linear-algebra operations.
\end{proof}

\begin{remark}[Rank zero and rejected cuts]
The explicit zero-flux equations on rank-zero cuts are necessary. A zero
coefficient $w_v$ alone would leave an unrelated auxiliary scalar free. A cut
with $\lambda_v>1$ is rejected, not projected onto a selected direction. A
forbidden pair crossing bags is likewise rejected, not silently dropped.
Both rejection paths are exercised by the tests.
\end{remark}

\subsection{An exact and reusable source contract}

The compiler accepts $A$, its declared number of columns, a partition into bags,
a rooted tree, and the complete forbidden-pair list. The returned model retains
original column names and the primitive cut generators. The source verifier
uses rank-and-membership checks, not the producer's intersection-basis routine.

If a source column is later disabled, the same compiled system with that
original coordinate set to zero remains exact. Rebuilding smaller cut spaces
may improve performance, but is not required for correctness. If the source
matching equations themselves change, such as after a Pachner move or crushing,
a new source compilation or a separately proved transport certificate is
required. An old small-width tree has no automatic authority over the new
triangulation.

\section{Source-bound gluing preserves extreme rays}
\label{sec:extreme}

An arbitrary signed scalar presentation can introduce artificial extreme rays.
For example, present $\R_+^2$ using the redundant free scalar $\tau=a-b$.
The chamber $\tau\geq0$ has the ray $(a,b)=(1,1)$, which is not extreme in
$\R_+^2$. It is therefore not enough to observe that a local signed ray was
chosen. The source-cut structure of Theorem~\ref{thm:compiler} repairs this gap.

For each sign vector $\sigma$ on the nonzero rank-one cuts, let
\[
 C_\sigma=\{x\in C(A):\sigma_v\tau_v(x)\geq0\text{ for every such }v\}.
\]
The scalars are now the unique source-bound linear functionals.

\begin{lemma}[Zero flux splits the source vector]
\label{lem:zerosplit}
If $x\in C(A)$ and $\tau_v(x)=0$, the restrictions of $x$ to $S_v$ and to
$\overline{S_v}$, extended by zero, both belong to $C(A)$. More generally,
remove every zero-flux edge from the tree. Restricting $x$ to the original
columns of each remaining tree component gives feasible vectors whose sum
is $x$. Each restricted vector belongs to every sign chamber containing $x$.
\end{lemma}

\begin{proof}
Equation~\eqref{eq:flux} gives $A_{S_v}x_{S_v}=0$, and the complementary
identity follows from $Ax=0$. For several edges, use the compiled equations:
retain a component's original coordinates and its internal fluxes and set all
other coordinates and fluxes to zero. Its boundary fluxes were already zero,
so every local equation remains true. The compiler theorem proves source
feasibility and identifies the retained fluxes as the unique source fluxes.
An internal flux keeps its sign, and an external one becomes zero.
\end{proof}

\begin{theorem}[No new source rays under flux-sign subdivision]
\label{thm:nonew}
Every extreme ray of every $C_\sigma$ is an extreme ray of $C(A)$. Conversely,
every extreme ray of $C(A)$ is an extreme ray of any $C_\sigma$ containing it.
Thus the union of the flux-sign chamber ray sets is exactly the source ray set.
The same statements hold after any original coordinate-face restrictions.
\end{theorem}

\begin{proof}
Let $x\ne0$ span an extreme ray of $C_\sigma$. By
Lemma~\ref{lem:zerosplit}, splitting the tree at the zero-flux edges decomposes
$x$ into feasible vectors in the same chamber with disjoint original supports.
Extremality implies that only one tree component carries a nonzero vector.
The source compiler's injectivity rules out a different component supported
only by nonzero auxiliary fluxes.

Let $F$ be the coordinate face of $C(A)$ obtained by setting every coordinate
outside $\supp x$ to zero. The vector $x$ lies in the relative interior of $F$:
it is positive on all allowed support coordinates, so sufficiently small
perturbations in $\operatorname{span}F$ preserve nonnegativity.

Every zero-flux edge separates the single occupied tree component entirely
onto one side. For any $y\in F$, all its coordinates are on that side, and
$Ay=0$ forces the cut flux of $y$ to be zero. Thus every sign inequality that
is an equality at $x$ is identically zero on $F$. All other sign inequalities
are strict at $x$ and remain valid under sufficiently small perturbations.
If $\dim F\geq2$, there is a vector $u\in\operatorname{span}F$ not parallel to
$x$ for which $x\pm\varepsilon u\in C_\sigma$ for small $\varepsilon>0$.
This contradicts extremality. Hence $F$ has dimension one, so $x$ is an extreme
ray of $C(A)$.

The converse is immediate: a decomposition in the smaller cone would be a
decomposition in the larger cone. Coordinate-zero restrictions commute with
all the restrictions and zero extensions in the argument, proving the last
assertion.
\end{proof}

\begin{theorem}[Ray gluing and connected primitive witnesses]
\label{thm:gluing}
In the source-compiled model, suppose the local-ray recurrence chooses a single
local extreme ray at each nonzero recipe node and uses the zero assignment on
zero-port child subtrees. Every resulting nonzero global witness spans an
extreme ray of the original matching cone, and satisfies the original
compatibility constraints. If the source is a genuine standard normal-surface
system, its primitive integer normalization represents a connected normal
surface.
\end{theorem}

\begin{proof}
Choose the global flux-sign chamber containing the reconstructed assignment.
Consider a decomposition of this assignment into two vectors of that chamber,
using the nonnegative variables obtained after fixing signs. At the origin
node, the local vector is an extreme ray, so the two local summands are scalar
multiples of it. Every nonzero child port forces the same scalar ratio at the
child: its chosen local ray is extreme and its parent port is nonzero. This
ratio propagates through all nonzero recipe edges. On a zero child subtree,
all nonnegative chamber coordinates are zero, so both summands vanish there.
Outside an early positive witness's origin subtree the same argument applies.
Thus the full assignment is extreme in the compiled chamber. Projection is
injective, and Theorem~\ref{thm:nonew} makes the source vector extreme in $C(A)$.
Original incompatibilities hold because they are local and every chosen ray
satisfies them.

Scale the source vector to the primitive nonnegative integer vector $p$ on its
ray. Normal matching and compatibility give an embedded normal surface with
coordinate vector $p$~\cite{haken}. If it had two nonempty connected components,
their coordinate vectors would be nonzero nonnegative integer vectors summing
to $p$. Extremality makes each a nonnegative scalar multiple of $p$. An integer
vector parallel to a primitive integer vector is an integer multiple: apply a
Bezout combination of the coordinates of $p$, whose gcd is one. Each nonzero
component would therefore contribute at least $p$, a contradiction.
\end{proof}

\begin{remark}[Why the stronger conclusion is source-sensitive]
The generic scalar solver is correct even for a model with redundant ports.
Its positive output need not then be source-extreme. The code includes the
$(a,b,a-b)$ counterexample and offers greedy original-support minimization as
a safe fallback for a genuinely source-equivalent generic model. For a model
accepted by the rank-one source compiler, the preceding theorem makes that
extra search unnecessary. Independent rank tests on the positive support
provide an inexpensive additional audit when the source matrix is available.
\end{remark}
